% =====================================================================
%  Appendix
% =====================================================================
\appendix

% =====================================================================
\section{Proof of Theorem~\ref{thm:stability}}
\label{app:proof-stability}

This appendix proves the per-epoch and cumulative bounds of
Theorem~\ref{thm:stability}, together with the score-transfer
corollary used in Sec.~\ref{sec:method-core}. The proof has three
steps: a Davis--Kahan angular bound, a conversion to Euclidean
distance via sign calibration, and a transfer from anchor stability
to the \textsc{TADS} score.

\paragraph{Assumption interpretations.}
Assumption~1 holds when parameter updates are learning-rate bounded
and the map $\theta \mapsto h_l^{\mathrm{last}}(\cdot;\theta)$ is
locally Lipschitz. Re-sampling the probe can add an
$\eta_t$-independent term of order
$O(|\widetilde{\mathcal{D}}|^{-1/2})$; empirically this term is small
at the probe sizes used in our verification. Assumption~2
($\gamma^{(t)}\ge\gamma_{\min}$) is checked in
App.~\ref{app:thm-verification}. Assumption~3 is enforced by sign
calibration in Alg.~\ref{alg:tads}.

\subsection{Step 1: Angular bound}
\label{app:step-angular}

The Davis--Kahan $\sin\Theta$ theorem~\cite{daviskahan} states that
for symmetric matrices $A,B$ with top unit eigenvectors $u,v$ and
top eigengap $\delta=\lambda_1(A)-\lambda_2(A)>0$,
\[
\sin\Theta(u,v)\le \frac{\|A-B\|_F}{\delta},
\qquad
\Theta(u,v)=\arccos|\langle u,v\rangle|.
\]
Apply this with $A=\Sigma^{(t)}$, $B=\Sigma^{(t+1)}$,
$\delta=\gamma^{(t)}\ge\gamma_{\min}$, and
$\|\Sigma^{(t+1)}-\Sigma^{(t)}\|_F\le C_\Sigma\eta_t$. Then
\begin{equation}
  \sin\Theta\bigl(v_l^{(t+1)},v_l^{(t)}\bigr)
  \le
  \frac{C_\Sigma}{\gamma_{\min}}\eta_t .
  \label{eq:angular-bound}
\end{equation}

\subsection{Step 2: Per-epoch and cumulative bounds}
\label{app:step-euclidean}

The angular bound is invariant to the PCA sign ambiguity. Sign
calibration ensures
$\langle v_l^{(t+1)},v_l^{(t)}\rangle>0$, so
$\Theta\in[0,\pi/2)$. Hence
\begin{equation}
  \bigl\|v_l^{(t+1)}-v_l^{(t)}\bigr\|
  =
  2\sin(\Theta/2)
  \le
  2\sin\Theta
  \le
  \frac{2C_\Sigma}{\gamma_{\min}}\eta_t .
  \label{eq:perepoch-bound}
\end{equation}
This is the per-epoch bound in Theorem~\ref{thm:stability}. Summing
over epochs gives
\begin{equation}
  \bigl\|v_l^{(t+m)}-v_l^{(t)}\bigr\|
  \le
  \frac{2C_\Sigma}{\gamma_{\min}}
  \sum_{s=t}^{t+m-1}\eta_s .
  \label{eq:cumulative-bound-app}
\end{equation}
Under a decaying learning-rate schedule, the tail
$\sum_{s\ge t}\eta_s$ shrinks with $t$, yielding the stated
$\varepsilon$-stationarity.

\subsection{Step 3: Score stability transfer}
\label{app:score-stability}

Write the \textsc{TADS} score as
\[
s_i^{(t)} = B_i^{(t)}A_i^{(t)},
\qquad
B_i^{(t)} := \mathcal{R}_i^{(t)} \in (0,1),
\qquad
A_i^{(t)} := 1+\lambda\widetilde{\mathrm{align}}_i^{(t)} .
\]
Here $B_i^{(t)}$ is the composite-reward base ranking
(Eq.~\ref{eq:composite}) and $A_i^{(t)}\in[1,1+\lambda]$ is the
anchor factor.

\begin{corollary}[Score stability transfer]
\label{cor:score-stability}
Assume $|B_i^{(t)}|\le B_{\max}$, the sequence-mean activations
satisfy $\|\bar h_l(x_i;\theta_{t-1})\|\le H_{\max}$, and the
min--max normalisation in Eq.~\ref{eq:align} has Lipschitz constant
$L_{\mathrm{norm}}$. Then the anchor-induced score change satisfies
\begin{equation}
  \bigl|
  B_i^{(t+1)}
  \bigl(A_i^{(t+1)}-A_i^{(t)}\bigr)
  \bigr|
  \le
  B_{\max}\lambda L_{\mathrm{norm}}H_{\max}
  \frac{2C_\Sigma}{\gamma_{\min}}\eta_t
  =
  O(\lambda\eta_t).
  \label{eq:cor-score-bound}
\end{equation}
\end{corollary}

\begin{proof}
By Cauchy--Schwarz,
$\langle \bar h_l(x_i;\theta_{t-1}),v_l^{(t)}\rangle$ is
$H_{\max}$-Lipschitz in $v_l^{(t)}$. Therefore the raw alignment
changes by at most
$H_{\max}\|v_l^{(t+1)}-v_l^{(t)}\|$. Composing with the
$L_{\mathrm{norm}}$-Lipschitz min--max normalisation and applying
Eq.~\ref{eq:perepoch-bound} gives
\[
|A_i^{(t+1)}-A_i^{(t)}|
\le
\lambda L_{\mathrm{norm}}H_{\max}
\frac{2C_\Sigma}{\gamma_{\min}}\eta_t .
\]
Multiplying by $|B_i^{(t+1)}|\le B_{\max}$ proves
Eq.~\ref{eq:cor-score-bound}.
\end{proof}

The corollary holds for any uniformly bounded base ranking score
$B_i^{(t)}$. For the composite reward of Eq.~\ref{eq:composite}
the boundedness assumption reduces to assuming bounded
$\mathcal{L}_i, \mathcal{H}_i$, which holds under the standard
finite-vocabulary, finite-sequence-length setup
($\mathcal{L}_i \le \log|V|\cdot K_{\max}$,
$\mathcal{H}_i \le \log|V|$). Thus the stability result concerns
the trajectory anchor and its multiplicative factor; it is not
tied to the particular form of the per-candidate signal used on
the left of Eq.~\ref{eq:score}.

% =====================================================================
\section{Experimental Setup Details}
\label{app:setup-details}

\paragraph{Hardware.}
NVIDIA A-series 80GB GPUs with NVLink. Main LLaMA-2-7B experiments:
single 8$\times$A100 node. Qwen2.5-14B LoRA scaling: $8\times$A100.

\paragraph{Main experiment: LLaMA-2-7B full fine-tuning.}
Adam ($\beta_1=0.9, \beta_2=0.999$), learning rate $2\times10^{-5}$,
cosine decay to zero, warmup ratio 0.03, weight decay 0.1, bf16
precision, 3 epochs, batch size 128, cutoff length 2{,}048,
gradient checkpointing, 8-bit AdamW (bitsandbytes). The
configuration is held fixed across selection methods.

\paragraph{Evaluation protocol.}
We evaluate with lm-evaluation-harness~\cite{eval-harness} on nine
tasks: MMLU, MMLU-Pro, GSM8K, SVAMP, HumanEval, MBPP, TydiQA,
XQuAD, and BBH. Per-task scores and AVG are reported. Concrete
prompts and scoring routines for each benchmark appear in
App.~\ref{app:examples}.

\paragraph{\textsc{TADS}-specific hyperparameters.}
Probe size $|\widetilde{\mathcal{D}}_t|=1{,}024$; trajectory anchor
refreshed each epoch; alignment averaged over all decoder layers;
default $\lambda=1.0$. Numerical stabiliser $\varepsilon=10^{-6}$
in Eq.~\ref{eq:w}.

\paragraph{Composite reward.}
At each refresh step we form
$\mathcal{R}_i^{(t)} = w^{(t)}\mathcal{L}_i^{(t)} +
(1-w^{(t)})\mathcal{H}_i^{(t)}$
with $w^{(t)}$ from Eq.~\ref{eq:w} computed pool-wide. The composite
reward is used as-is in Eq.~\ref{eq:score}; no z-score, sigmoid,
clipping, or learned transform is applied between $\mathcal{R}_i^{(t)}$
and the ranking step.

\paragraph{Implementation note.}
Our public implementation includes an optional PPO actor--critic
substrate retained for compatibility with future extensions to
verifiable-reward settings (RLVR/GRPO). The experiments in this
paper do not consume the actor's output in the ranking rule; all
reported numbers use the composite-reward ranking of
Eq.~\ref{eq:score} directly.

% =====================================================================
\section{$\lambda$-Ablation and Qwen2.5-0.5B Sanity Track}
\label{app:setup-validation}

Before the main experiments, we validate the pipeline on
Qwen2.5-0.5B with full fine-tuning on Alpaca-GPT4 ($52{,}002$ pairs),
$\rho=10\%$, 3 epochs, AdamW, lr $2\times10^{-5}$, cosine schedule
with $10\%$ warmup, MMLU + GSM8K evaluation.

The $\lambda$-ablation (Table~\ref{tab:lambda-05b}) is
non-monotonic: small non-zero $\lambda$ improves over the
$\lambda=0$ composite-reward base, while large $\lambda$
saturates.

% =====================================================================
\section{Scaling Study: Qwen2.5-14B}
\label{app:scaling}

Qwen2.5-14B + LoRA ($r=16, \alpha=32$, dropout $0.05$, target
modules $\{q,k,v,o,\text{gate},\text{up},\text{down}\}$proj),
Alpaca-GPT4 ($N=52{,}002$), 3 epochs, batch size 64, lr
$2\times10^{-4}$, cosine, $3\%$ warmup, $8\times$A100. Cost
breakdown in Table~\ref{tab:cost-14b-main}; \textsc{TADS}-only
operations (anchor PCA + alignment scoring) add $\sim$15s to a
$\sim$35-minute epoch, i.e. $<0.5\%$ overhead.

% =====================================================================
\section{Empirical Verification of Theorem~\ref{thm:stability}}
\label{app:thm-verification}

We instrument a Qwen2.5-0.5B LoRA run on Alpaca-GPT4 with a
verifier every 25 optimiser steps on a fixed 2{,}000-sample probe.
For each of 24 decoder layers, the verifier computes
$\gamma_l^{(t)}, \|\Sigma_l^{(t+1)}-\Sigma_l^{(t)}\|_F,$ and
$d_l^{(t)}=\|v_l^{(t+1)}-v_l^{(t)}\|$.

\paragraph{Assumptions hold.}
$\gamma_{\min}=0.050$; trimmed $\hat C_\Sigma=5.78\times10^3$;
sign-flip rate $3.6\%$ (below $5\%$ tolerance). The empirical
per-step bound at $\eta=2\times10^{-4}$ is
$(2\hat C_\Sigma / \gamma_{\min}) \eta = 46.24$, never violated
across $1{,}496$ measurements; median tightness
$1.46\times10^{-3}$, worst $7.4\times10^{-3}$.

\paragraph{No-anchor diagnostic.}
The matched no-anchor control (\texttt{use\_anchor=false},
$\lambda=0$) yields $2.9\times$ larger median displacement, a
single-layer maximum of $1.998$, and sign-flip rate above $5\%$
(Table~\ref{tab:dstep-spread}).

% =====================================================================
\section{In-Progress Empirical Study}
\label{app:inprogress}

The broader study extends along three axes: (i) scaling across
model sizes (0.5B / 7B / 14B); (ii) robustness across additional
IT corpora beyond Alpaca-GPT4 (Evol-Instruct-style); (iii)
ablations over $\rho$, anchor refresh frequency, layer subsets,
sign calibration, probe size.

% =====================================================================
\section{Dataset and Benchmark Examples}
\label{app:examples}

This appendix provides one concrete worked example of the candidate
pool (Alpaca-GPT4) and of each of the nine evaluation benchmarks,
to make the scoring conventions explicit.

\subsection{Candidate pool — Alpaca-GPT4}
\label{app:ex-alpaca}

Each candidate $x_i \in \mathcal{D}$ is a JSON object with three
fields: \texttt{instruction}, \texttt{input} (optional, present in
${\sim}40\%$ of samples), and \texttt{output}.

\paragraph{Example (one item from Alpaca-GPT4).}
\begin{verbatim}
{
  "instruction": "Translate the following English sentence to French.",
  "input":       "Hello, how are you doing today?",
  "output":      "Bonjour, comment allez-vous aujourd'hui ?"
}
\end{verbatim}

\paragraph{Tokenisation (Alpaca default template).}
The fields are flattened into
\verb|### Instruction:\n{instruction}\n\n### Input:\n{input}\n\n### Response:\n|
followed by the \texttt{output}. The template and instruction
tokens are masked with $-100$ in the labels; only response tokens
contribute to the cross-entropy loss used in $\mathcal{L}_i$.

\subsection{MMLU}
\label{app:ex-mmlu}

14{,}042 four-choice items across 57 subjects. Five-shot prompting;
scoring by next-token logit comparison.

\paragraph{Example item (subject: high\_school\_biology).}
\begin{verbatim}
Q: Which of the following organelles is responsible for
   producing most of a eukaryotic cell's ATP?
(A) Ribosome   (B) Mitochondrion
(C) Nucleus    (D) Golgi apparatus
Answer:
\end{verbatim}

\paragraph{How a model "matches".} Predict the next token after
\verb|Answer:|. We extract logits at four candidate-letter token
positions and pick the largest. Correct iff the chosen letter is
\verb|B|.

\subsection{MMLU-Pro}
\label{app:ex-mmlu-pro}

12{,}032 items, ten options per question. Five-shot CoT.

\paragraph{Example item (subject: physics).}
\begin{verbatim}
Q: A particle moves in a circle of radius 2 m at constant
   speed 3 m/s. What is the magnitude of its centripetal
   acceleration?
Options: (A) 0.67 (B) 1.5 (C) 3.0 (D) 4.5 (E) 6.0
         (F) 9.0 (G) 12.0 (H) 18.0 (I) 24.0 (J) 36.0
Answer: Let's think step by step.
\end{verbatim}

\paragraph{How a model "matches".} Model generates a CoT ending in
\verb|The answer is (X)|. Extracted letter must equal the gold
label (\verb|D|: $a=v^2/r=9/2=4.5$).

\subsection{GSM8K}
\label{app:ex-gsm8k}

1{,}319 grade-school word problems. Eight-shot CoT; exact-match on
the final numeric answer.

\paragraph{Example item.}
\begin{verbatim}
Q: Natalia sold clips to 48 of her friends in April, and then
   she sold half as many clips in May. How many clips did
   Natalia sell altogether in April and May?
A: Let's think step by step. In April she sold 48. In May she
   sold 48 / 2 = 24. Altogether she sold 48 + 24 = 72.
   The answer is 72.
\end{verbatim}

\paragraph{How a model "matches".} Extract the integer after
\verb|The answer is| (fallback: last numeric in text). Equality to
gold target $72$ counts as correct.

\subsection{SVAMP}
\label{app:ex-svamp}

1{,}000 arithmetic word problems with adversarial phrasings.
Eight-shot CoT.

\paragraph{Example item.}
\begin{verbatim}
Body: Each child has 3 pencils and 4 erasers. There are 8 children.
Q: How many pencils are there in total?
A: Each child has 3 pencils, and there are 8 children.
   Total = 3 * 8 = 24. The answer is 24.
\end{verbatim}

\paragraph{How a model "matches".} Same as GSM8K — exact match on
trailing integer.

\subsection{HumanEval}
\label{app:ex-humaneval}

164 hand-written Python programming problems with hidden unit
tests. Zero-shot; pass@1 by execution.

\paragraph{Example item.}
\begin{verbatim}
def has_close_elements(numbers: List[float],
                       threshold: float) -> bool:
    """ Check if in given list of numbers, are any two numbers
    closer to each other than given threshold.
    >>> has_close_elements([1.0, 2.0, 3.0], 0.5)
    False
    >>> has_close_elements([1.0, 2.8, 3.0, 4.0, 5.0, 2.0], 0.3)
    True
    """
\end{verbatim}

\paragraph{How a model "matches".} Model greedily generates the
function body. We concatenate prompt + completion, write to a
temp file, and run hidden unit tests in a sandboxed subprocess
with 3s timeout. The problem is passed iff all assertions hold;
pass@1 is the fraction of problems passed.

\subsection{MBPP}
\label{app:ex-mbpp}

500 test-split Python problems with three given unit tests per
problem. Three-shot; pass@1 by execution.

\paragraph{Example item.}
\begin{verbatim}
"""
Write a function to find the minimum cost path to reach (m, n)
from (0, 0) for the given cost matrix cost[][].
"""
TESTS:
assert min_cost([[1,2,3],[4,8,2],[1,5,3]], 2, 2) == 8
assert min_cost([[2,3,4],[5,9,3],[2,6,4]], 2, 2) == 12
assert min_cost([[3,4,5],[6,10,4],[3,7,5]], 2, 2) == 16
\end{verbatim}

\paragraph{How a model "matches".} Concatenate the completion with
the three asserts; if all asserts pass without raising, the problem
is passed. pass@1 is the fraction passed.

\subsection{TydiQA}
\label{app:ex-tydiqa}

Multilingual gold-passage QA in nine languages (en, ar, bn, fi, id,
ko, ru, sw, te). One-shot; token-overlap F1.

\paragraph{Example item (Korean).}
\begin{verbatim}
Passage: 에펠탑은 1889년 파리 만국박람회를 위해 Gustave Eiffel
회사가 건설한 격자형 철탑이다. 높이는 약 330m이며, 건설 당시
세계에서 가장 높은 인공 구조물이었다.
Q: 에펠탑의 높이는 얼마인가?
A: 약 330m
\end{verbatim}

\paragraph{How a model "matches".} After lowercasing, removing
articles, and stripping punctuation, we compute token precision and
recall between the predicted answer and each gold answer, taking
the maximum F1. Per-language F1 is averaged over the nine
languages.

\subsection{XQuAD}
\label{app:ex-xquad}

1{,}190 parallel SQuAD-style items in eleven languages. One-shot;
SQuAD-style token-overlap F1.

\paragraph{Example item (Spanish).}
\begin{verbatim}
Passage: Los Beatles fueron una banda de rock inglesa formada
en Liverpool en 1960. Estaba integrada por John Lennon, Paul
McCartney, George Harrison y Ringo Starr.
Q: ¿En qué ciudad se formaron los Beatles?
A: Liverpool
\end{verbatim}

\paragraph{How a model "matches".} Same F1 procedure as TydiQA;
report average over eleven languages.

\subsection{BBH}
\label{app:ex-bbh}

Big-Bench Hard: 27 sub-tasks. Three-shot CoT.

\paragraph{Example item (sub-task: logical\_deduction\_three\_objects).}
\begin{verbatim}
Q: Three friends, Alice, Bob, and Carol, finished a race.
   Alice finished after Bob. Carol finished before Bob.
   Who finished last?
   Options: (A) Alice  (B) Bob  (C) Carol
A: Let's think step by step. Carol < Bob < Alice, so Alice
   finished last. The answer is (A).
\end{verbatim}

\paragraph{How a model "matches".} Model ends with
\verb|The answer is (X)| or with a numeric/textual answer. The
per-sub-task scorer compares against the gold; BBH score is the
unweighted mean accuracy across 27 sub-tasks.
